% Additive sharpening note for the D0033 source-scope package.
% R7-01, 2026-09-17. This file does not replace main.tex.

\section{A Dobrushin--SDPI sharpening of the quantitative bound}
\label{sec:sdpi-sharpening}

This section records a quantitative strengthening of the source-level argument.
The source question only asks whether $\phi(n)\to0$ forces $\rho^*(n)\to0$ for a
Markov chain.  The earlier proof gave the sufficient estimate
$\rho^*(n)\le 4\sqrt{2\phi(n)}$.  The local step can be sharpened by changing
language from a direct $\phi$--to--$\rho$ inequality to a channel contraction
inequality.

For a Markov transition kernel $K$ define its Dobrushin row diameter by
\[
  \delta(K):=\operatorname*{ess\,sup}_{x,x'}
  \|K(x,\cdot)-K(x',\cdot)\|_{\rm TV},
\]
where the essential support is taken with respect to the relevant input law.
For the time-inhomogeneous Markov chain $X=(X_k)$ write
\[
  \Delta(n):=\sup_{s<t:\,t-s\ge n}\delta(K_{s,t}).
\]
The convention for total variation is
$\|\mu-\nu\|_{\rm TV}=\sup_B|\mu(B)-\nu(B)|$.

\begin{lemma}[Dobrushin diameter controls maximal correlation]
\label{lem:dobrushin-rho}
Let $(U,V)$ be random variables on standard Borel spaces and let $K(u,\cdot)$
be a regular conditional law of $V$ given $U=u$.  If the essential row
diameter of $K$ is at most $\delta$, then
\[
  \rho(U,V)^2\le \delta.
\]
The same conclusion holds after replacing the law of $U$ by any probability
measure supported on the same row-good set.
\end{lemma}

\begin{proof}
We first treat finite-valued $U,V$.  For a finite stochastic matrix $K$, the
classical contraction theorem of Cohen--Iwasa--Rautu--Ruskai--Seneta--Zbaganu
for relative $g$-entropies states that every $f$-divergence contraction
coefficient is at most the Dobrushin coefficient.  Taking
$f(t)=(t-1)^2$ gives
\[
  \eta_{\chi^2}(K)\le\delta(K).
\]
For completeness, the link with maximal correlation is immediate.  Let
$\lambda$ be the law of $U$, let $h\in L^2_0(\lambda)$ be bounded, and for
$|\varepsilon|$ small enough put
\[
  d\lambda_\varepsilon=(1+\varepsilon h)\,d\lambda.
\]
If $\nu=\lambda K$, then
\[
 \frac{d(\lambda_\varepsilon K)}{d\nu}(v)
   =1+\varepsilon\,\mathbb E[h(U)\mid V=v].
\]
Consequently
\[
 \frac{\chi^2(\lambda_\varepsilon K\|\nu)}
      {\chi^2(\lambda_\varepsilon\|\lambda)}
 =\frac{\|\mathbb E[h(U)\mid V]\|_2^2}{\|h\|_2^2}.
\]
Taking the supremum over bounded centered $h$, which are dense in
$L^2_0(\lambda)$, gives
\[
  \rho(U,V)^2\le \eta_{\chi^2}(K)\le\delta(K).
\]

For standard Borel $U,V$, approximate arbitrary centered square-integrable
functions of $U$ and $V$ in $L^2$ by bounded simple functions.  The induced
finite channel has rows that are push-forwards of convex mixtures of rows of
$K$.  If all original rows have pairwise total-variation distance at most
$\delta$, then any two convex mixtures of those rows also have distance at
most $\delta$, because
\[
 \Bigl\|\int K(u,\cdot)\,d\alpha(u)-
        \int K(u',\cdot)\,d\beta(u')\Bigr\|_{\rm TV}
 \le \iint\|K(u,\cdot)-K(u',\cdot)\|_{\rm TV}
          \,d\alpha(u)d\beta(u')
 \le\delta.
\]
Push-forward under a finite partition can only decrease total variation.
The finite result therefore bounds the correlation of every pair of simple
approximants by $\sqrt\delta$, and passage to the $L^2$ limit proves the
claim.  The proof depends only on the row family, so changing the input law
while remaining on the row-good set does not change the estimate.
\end{proof}

\begin{theorem}[Intrinsic row-diameter bound]
\label{thm:row-diameter-rhostar}
For a real-state Markov chain in Bradley's Section~3 setting,
\[
  \rho^*(n)\le \min\{1,\,2\sqrt{\Delta(n)}\}.
\]
In particular, Bradley's $\phi$ coefficient satisfies
\[
  \Delta(n)\le \min\{1,2\phi(n)\},
\]
and hence
\[
  \boxed{\rho^*(n)\le \min\{1,\,2\sqrt{2\phi(n)}\}.}
\]
Thus $\phi$-mixing implies $\rho^*$-mixing.
\end{theorem}

\begin{proof}
The earlier row-diameter lemma in the source-scope proof shows that for each
$s<t$, outside one common null set,
\[
  \|K_{s,t}(x,\cdot)-K_{s,t}(x',\cdot)\|_{\rm TV}
  \le 2\phi(t-s).
\]
This proves $\Delta(n)\le\min\{1,2\phi(n)\}$.

Now take finite disjoint index sets $S,T$ with $\operatorname{dist}(S,T)\ge n$.
Order $S\cup T$ chronologically and merge consecutive selected indices of the
same colour into maximal monochromatic blocks.  The block colours alternate.
Enumerate the $S$-blocks as $S_1,\dots,S_I$ and the $T$-blocks as
$T_1,\dots,T_J$, and let $U_i=(X_k)_{k\in S_i}$ and
$V_j=(X_k)_{k\in T_j}$.

Use Peyre's minimal-hypothesis form of the $N$-against-$M$ tensorization
theorem: for the pair $(U_i,V_j)$ one conditions only on the earlier
$S$-blocks and earlier $T$-blocks.  If $U_i$ and $V_j$ are nonadjacent in the
full chronological block path, that conditioning field contains an
intervening block and hence an intervening Markov state.  The Markov property
then gives subjective maximal correlation exactly zero.

If they are adjacent, write $s$ for the last time of the earlier block and
$t$ for the first time of the later block.  Then $t-s\ge n$.  Conditioning on
the prescribed earlier blocks changes the input law of $X_s$ but not the
transition kernel $K_{s,t}$; the conditional input law remains supported on
the row-good set.  Lemma~\ref{lem:dobrushin-rho} and maximal-correlation data
processing give
\[
  \rho(U_i,V_j\mid\hbox{earlier blocks})
  \le \sqrt{\delta(K_{s,t})}
  \le \sqrt{\Delta(n)}.
\]
Thus Peyre's subjective-correlation matrix is the bipartite adjacency matrix
of a path with every nonzero entry at most $\sqrt{\Delta(n)}$.  Its
$\ell^2\to\ell^2$ norm is at most
\[
  \sqrt{\|E\|_1\|E\|_\infty}\le 2\sqrt{\Delta(n)},
\]
because every row and every column has at most two nonzero entries.  Peyre's
tensorization theorem yields the same bound for the maximal correlation of
the two finite coordinate vectors.  Arbitrary index sets follow by the usual
finite-coordinate $L^2$ approximation of their generated sigma-fields.
\end{proof}

\paragraph{Sharpness of the local square root.}
The exponent in Lemma~\ref{lem:dobrushin-rho} cannot be improved from
$\sqrt\delta$ using Dobrushin diameter alone.  For the binary erasure channel
which reveals the input with probability $\delta$ and otherwise outputs an
erasure symbol, the two channel rows have total-variation distance $\delta$.
With a symmetric Bernoulli input, conditional expectation preserves exactly
a $\delta$ fraction of every centered input variance, so
\[
  \rho(U,V)^2=\delta.
\]
Therefore any further improvement to the global constant must use additional
Markov/stationarity structure or improve the path tensorization step; it
cannot come from a stronger universal conversion of row diameter to local
maximal correlation.

\paragraph{Prior-art note.}
The finite-channel inequality used above is classical: J.~E. Cohen,
Y.~Iwasa, Gh.~Rautu, M.~B. Ruskai, E.~Seneta, and Gh.~Zbaganu,
``Relative entropy under mappings by stochastic matrices,'' Linear Algebra
Appl. 179 (1993), 211--235.  Modern strong-data-processing treatments include
Y.~Polyanskiy and Y.~Wu, arXiv:1508.06025, and O.~Ordentlich and
Y.~Polyanskiy, IEEE Trans. Inf. Theory 68 (2022), 1480--1481.  The global
block step uses R.~Peyre, ``Tensorizing maximal correlations,''
arXiv:1004.1602v2 (2010), Theorem~3.3.1 and Section~3.4(a).  These ingredients
are prior art.  The assertion that their combination resolves Bradley's
Question~5 remains publication-novelty UNVERIFIED pending a broader literature
review and the required independent audits.
